\scanpage{186}
désignant la ponctuation. Plus précisément, si on a un autre foncteur $\frac12$\inserted{co}simplicial
\diagram{\begin{tikzcd}[ampersand replacement=\&,column sep=1em]
f^0\arrow[r,shift left=.6ex]\arrow[r,shift right=.6ex]\&f^1\arrow[r,shift left=1.1ex]\arrow[r]\arrow[r,shift right=1.1ex]\&f^{\unsure{2|?}}
\end{tikzcd}}
\deleted{on va \unsure{se}} et un hom $\frac12$\inserted{co}simplicial $\xi$
\diagram{\begin{tikzcd}[ampersand replacement=\&,column sep=1em,row sep=1.4em]
\&F^0\arrow[r,shift left=.6ex]\arrow[r,shift right=.6ex]\arrow[d,"\xi^0"]\& F^1\arrow[r,shift left=1.1ex]\arrow[r]\arrow[r,shift right=1.1ex]\arrow[d,"\xi^1"]\&F^2\arrow[d,"\xi^2"]\\
\xi:\&f^0\arrow[r,shift left=.6ex]\arrow[r,shift right=.6ex]\&f^1\arrow[r,shift left=1.1ex]\arrow[r]\arrow[r,shift right=1.1ex]\&f^2
\end{tikzcd}}
on définit $\pi_0(F^*,\xi;E)$ et $\pi^{\unsure{2|1}}(F^*,\xi;E)$.

Si $F^0,F^1,F^2$ sont T-représentables, et dans le cas où on \unsure{se donne|les a un} \deleted{\unsure{de}} $\xi:F^*\unsure{\mathbin{-}|\to} f^*$, si $f^*$ est T-représentable, on ramène ceci à la donnée d'un \deleted{objet} \uline{espace $\frac12$simplicial}
\diagram{\begin{tikzcd}[ampersand replacement=\&,column sep=1em]
K_0\&K_1\arrow[l,shift left=.6ex]\arrow[l,shift right=.6ex]\&K_2\arrow[l,shift left=1.1ex]\arrow[l]\arrow[l,shift right=1.1ex]
\end{tikzcd}}
resp. d'un \uline{espace $\frac12$simplicial} \uuline{ponctué} (par contraction)
\diagram{\begin{tikzcd}[ampersand replacement=\&,column sep=1em]
\overline K_0\&\overline K_1\arrow[l,shift left=.6ex]\arrow[l,shift right=.6ex]\&\overline K_2\arrow[l,shift left=1.1ex]\arrow[l]\arrow[l,shift right=1.1ex]
\end{tikzcd}}
par un procédé \unsure{standard|standart} .....

Il y a lieu de généraliser ces constructions au cas où sur les $K_i$, resp. des $\overline K_i$, on a des catégories galoisiennes \unsure{attachées|étalées|itérées} à chaque pt, \deleted{(\unsure{ou encore}, les \unsure{pts|$F_\xi$}} pour décrire des relations de descente à la \uline{Van Kampen} .....

% archive sheet number: 185
